% =============================================================================
%  《量子场论与 Jones 多项式》读书笔记 —— 导言区
%  编译：xelatex（两遍）；中文字体：fandol
% =============================================================================
\documentclass[11pt,a4paper]{article}

\usepackage[UTF8, fontset=fandol, heading=true]{ctex}
\usepackage{amsmath}
\usepackage{amssymb}
\usepackage{mathtools}
\usepackage{bm}
\usepackage[a4paper,left=2.6cm,right=2.6cm,top=2.7cm,bottom=2.7cm]{geometry}
\usepackage{graphicx}
\usepackage{xcolor}
\usepackage{tcolorbox}
\tcbuselibrary{breakable, skins}
\usepackage{enumitem}
\usepackage{booktabs}
\usepackage{longtable}
\usepackage{array}
\usepackage[font=small,labelfont=bf,labelsep=quad]{caption}
\usepackage[colorlinks=true,
            linkcolor=blue!45!black,
            citecolor=blue!45!black,
            urlcolor=blue!45!black]{hyperref}

\linespread{1.15}
\setlength{\emergencystretch}{3em}
\allowdisplaybreaks[3]

\renewcommand{\figurename}{图}
\renewcommand{\tablename}{表}
\renewcommand{\refname}{参考文献}
\renewcommand{\contentsname}{目录}

\hypersetup{
  pdftitle={量子场论与 Jones 多项式：读书笔记},
  pdfauthor={WIT-1989-JONES 研讨班},
  pdfsubject={Chern-Simons theory, Jones polynomial, geometric quantization},
  pdfkeywords={Chern-Simons, Jones polynomial, conformal blocks, Verlinde formula},
  bookmarksnumbered=true
}

\newcommand{\ii}{\mathrm{i}}
\DeclareMathOperator{\Tr}{Tr}
\DeclareMathOperator{\Hom}{Hom}
\DeclareMathOperator{\Inv}{Inv}
\newcommand{\dd}{\mathrm{d}}
\newcommand{\CC}{\mathbb{C}}
\newcommand{\ZZ}{\mathbb{Z}}
\newcommand{\cM}{\mathcal{M}}
\newcommand{\cL}{\mathcal{L}}
\newcommand{\cH}{\mathcal{H}}
\newcommand{\cO}{\mathcal{O}}
\newcommand{\cD}{\mathcal{D}}
\newcommand{\cN}{\mathcal{N}}

% ---- 三种断言标记：原文 / 我们推导 / 我们猜测 -------------------------------
\newcommand{\src}[1]{\textnormal{\footnotesize[原文 p.#1]}}
\newcommand{\ours}{\textnormal{\footnotesize[我们推导]}}
\newcommand{\guess}{\textnormal{\footnotesize[我们猜测]}}

% ---- 章末「边界」框 ---------------------------------------------------------
\newenvironment{boundary}{%
    \par\medskip
    \begin{tcolorbox}[enhanced, breakable,
        colback=orange!4!white, colframe=orange!55!black,
        boxrule=0.5pt, arc=2pt, left=6pt, right=6pt, top=4pt, bottom=4pt,
        title={\small\textbf{本章边界}\quad{\footnotesize（哪些已证、哪些未证）}},
        fonttitle=\bfseries, coltitle=black, colbacktitle=orange!12!white]
}{%
    \end{tcolorbox}
    \par\medskip
}

% ---- 关键结论框 -------------------------------------------------------------
\newenvironment{keypoint}{%
    \par\medskip
    \begin{tcolorbox}[enhanced, breakable,
        colback=blue!3!white, colframe=blue!40!black,
        boxrule=0.5pt, arc=2pt, left=6pt, right=6pt, top=4pt, bottom=4pt]
}{%
    \end{tcolorbox}
    \par\medskip
}

% ---- 排版容差 ----
\tolerance=3000
\emergencystretch=3em
\hbadness=10000
